\subsection{\texorpdfstring{\(\log(1 + x)\)、Arc tan x 与 Arc sin x 的展开}{log(1 + x)、Arc tan x 与 Arc sin x 的展开}}

将 (26) 两端从 0 到 x 积分，即得 \(\log(1 + x)\) 的 n 阶泰勒展开，它对 x \textgreater{} -1 成立

\[
\log (1 + x) = \frac {x}{1} - \frac {x ^ {2}}{2} + \frac {x ^ {3}}{3} + \dots + (- 1) ^ {n - 1} \frac {x ^ {n}}{n} + (- 1) ^ {n} \int_ {0} ^ {x} \frac {t ^ {n} d t}{1 + t}.\tag{27}
\]

余项与 \((-1)^{n}\) 同号（当 x \textgreater{} 0 时），且为 \textless{} 0 （当 -1 \textless{} x \textless{} 0 时）；再者，当 x \textgreater{} 0 时，有 \(1 + t \geqslant 1\) （对 \(0 \leqslant t \leqslant x\) ），而当 -1 \textless{} x \textless{} 0 时，有 \(1 + t \geqslant 1 - |x|\) （对 \(x \leqslant 0\) ）；由此得余项的估计式

\[
\left| \int_ {0} ^ {x} \frac {t ^ {n} d t}{1 + t} \right| \leqslant \frac {| x | ^ {n + 1}}{n + 1}
\]

\[
\text { for } x \geqslant 0\tag{28}
\]

\[
\left| \int_ {0} ^ {x} \frac {t ^ {n} d t}{1 + t} \right| \leqslant \frac {| x | ^ {n + 1}}{(n + 1) (1 - | x |)} \quad \text {   for   } - 1 <   x \leqslant 0.\tag{29}
\]

由最后这两个公式立即可得：当 \(-1 < x \leqslant 1\) 时有

\[
\log (1 + x) = \sum_ {n = 1} ^ {\infty} (- 1) ^ {n - 1} \frac {x ^ {n}}{n}\tag{30}
\]

该级数在含于 \(]-1,+1]\) 的每个紧区间上一致收敛，且当 \(|x|<1\) 时绝对收敛。

另一方面，当 \(|x| > 1\) 时 (30) 右端级数的通项的绝对值随 n 无限增大（III, p.~106）。当 x = -1 时该级数化为调和级数，其和为 \(+\infty\) （Gen.~Top., IV, p.~365）。

类似地，在 (26) 中把 \(x\) 换成 \(x^2\) ，并将两端从 0 到 \(x\) 积分，即得 \(2n - 1\) 阶泰勒展开（对 \(\operatorname{Arctan} x\) ），它对一切实数 \(x\) 成立

\[
\operatorname{Arc} \tan x = \frac {x}{1} - \frac {x ^ {3}}{3} + \frac {x ^ {5}}{5} + \dots + (- 1) ^ {n - 1} \frac {x ^ {2 n - 1}}{2 n - 1} + (- 1) ^ {n} \int_ {0} ^ {x} \frac {t ^ {2 n} d t}{1 + t ^ {2}}.\tag{31}
\]

余项与 \((-1)^n x\) 同号；由于 \(1 + t^2 \geqslant 1\) 对一切 \(t\) 成立，我们有估计式

\[
\left| \int_ {0} ^ {x} \frac {t ^ {2 n} d t}{1 + t ^ {2}} \right| \leqslant \frac {| x | ^ {2 n + 1}}{2 n + 1}\tag{32}
\]

由此推出：当 \(|x| \leqslant 1\) 时

\[
\operatorname{Arc} \tan x = \sum_ {n = 1} ^ {\infty} (- 1) ^ {n - 1} \frac {x ^ {2 n - 1}}{2 n - 1}\tag{33}
\]

该级数在 \([-1,+1]\) 上一致收敛，且当 \(|x|<1\) 时绝对收敛。

特别地，取 \(x = 1\) 便得到公式

\[
{\frac {\pi}{4}} = 1 - {\frac {1}{3}} + {\frac {1}{5}} + \dots + (- 1) ^ {n} {\frac {1}{2 n + 1}} + \dots .\tag{34}
\]

当 \(|x| > 1\) 时 (33) 右端的通项的绝对值随 \(n\) 无限增大。

最后，关于 Arc sin x 的泰勒展开，我们从其导数 \((1 - x^{2})^{-1/2}\) 的展开出发；后一展开是在二项级数展开中以 \(-x^{2}\) 代换 x 而从 \((1 + x)^{-1/2}\) 得到的；对 \(|x| < 1\) 这给出

\[
(1 - x ^ {2}) ^ {- 1 / 2} = 1 + \frac {1}{2} x ^ {2} + \frac {1 . 3}{2 . 4} x ^ {4} + \dots + \frac {1 . 3 . 5 \dots (2 n - 1)}{2 . 4 . 6 \dots 2 n} x ^ {2 n} + r _ {n} (x)
\]

并且由 (23) ，其余项满足界

\[
0 \leqslant r _ {n} (x) \leqslant \frac {x ^ {2 n + 2}}{\sqrt {1 - x ^ {2}}}
\]

此即对余项的界。

取前一展开式的原函数，便得

\[
\operatorname{Arc} \sin x = x + \frac {1}{2} \frac {x ^ {3}}{3} + \frac {1 . 3}{2 . 4} \frac {x ^ {5}}{5} + \dots + \frac {1 . 3 . 5 \dots (2 n - 1)}{2 . 4 . 6 \dots 2 n} \frac {x ^ {2 n + 1}}{2 n + 1} + \mathrm{R} _ {n} (x)\tag{35}
\]

其中 \(\mathrm{R}_{n}(x)\) 与 x 同号且满足不等式

\[
| \mathrm{R} _ {n} (x) | \leqslant \int_ {0} ^ {x} \frac {t ^ {2 n + 2} d t}{\sqrt {1 - t ^ {2}}}.\tag{36}
\]

此外，关系式 (35) 表明 \(\mathbf{R}_n(x)\) 当 \(x\) 趋于 1 或 \(-1\) 时趋于一个极限，故有

\[
| \mathrm{R} _ {n} (1) | \leqslant \int_ {0} ^ {1} \frac {t ^ {2 n + 2} d t}{\sqrt {1 - t ^ {2}}}.\tag{37}
\]

但 (37) 的右端当 \(n\) 趋于 \(+\infty\) 时趋于 0 ：事实上，由于积分 \(\int_0^1 dt / \sqrt{1 - t^2}\) 收敛，对每个 \(\varepsilon > 0\) 存在 \(a\) 满足 \(0 < a < 1\) 且 \(\int_a^1 dt / \sqrt{1 - t^2} \leqslant \varepsilon\) ；另一方面我们有

\[
\int_ {0} ^ {a} \frac {t ^ {2 n + 2} d t}{\sqrt {1 - t ^ {2}}} \leqslant \frac {1}{\sqrt {1 - a ^ {2}}} \int_ {0} ^ {a} t ^ {2 n + 2} d t = \frac {a ^ {2 n + 3}}{(2 n + 3) \sqrt {1 - a ^ {2}}}
\]

于是存在 \(n_0\) ，使得当 \(n \geqslant n_0\) 时 \(\frac{a^{2n+3}}{(2n+3)\sqrt{1-a^2}} \leqslant \varepsilon\) ；由此，最终得 \(|\mathrm{R}_n(x)| \leqslant 2\varepsilon\) （当 \(|x| \leqslant 1\) 且 \(n \geqslant n_0\) ）。于是有

\[
\operatorname{Arc} \sin x = \sum_ {n = 0} ^ {\infty} \frac {1 . 3 . 5 \dots (2 n - 1)}{2 . 4 . 6 \dots 2 n} \frac {x ^ {2 n + 1}}{2 n + 1}\tag{38}
\]

其右端在紧区间 \([-1, 1]\) 上正规收敛。

在相反的情形，可以像对二项级数那样证明：当 \(|x| > 1\) 时 (38) 右端级数的通项的绝对值无限增大。

例如，在 (38) 中取 \(x = \frac{1}{2}\) ，就得到数 \(\pi\) 的一个新表达式；

\[
\frac {\pi}{6} = \sum_ {n = 0} ^ {\infty} \frac {1 . 3 . 5 \dots (2 n - 1)}{2 . 4 . 6 \dots 2 n} \frac {1}{(2 n + 1) 2 ^ {2 n + 1}}
\]

它比公式 (34) 更适合于计算 \(\pi\) 的近似值（参见 Calcul numérique）；由此得到

\[
\pi = 3. 1 4 1   5 9 2   6 5 3 \dots
\]

精确到 \(1/10^{9}\) 。\(^{3}\)
% semantic-fragments: legacy-input-seam
\relax{}
\setcounter{exercisegroup}{0}
\refstepcounter{exercisegroup}
